\begin{landscape}
\begin{table}[H]
\centering
\setlength{\tabcolsep}{2pt}
\caption[Strategy Logistic Regression Core]{Strategy Logistic Regression for packages against number of commits to predict Core developers}
\begin{tabular}{r l r r r r r r r r r}
\toprule
\makecell[t]{\textbf{Number of}\\\textbf{First}\\\textbf{Commits}} & \makecell[t]{\textbf{Model}\\\textbf{Parameter}} & \makecell[t]{\textbf{Number of}\\\textbf{Developers}} & \makecell[t]{\textbf{Not Core}\\\textbf{$\mu$}} & \textbf{Core $\mu$} & \textbf{Odds Ratio} & \textbf{Std. Err} & \textbf{$p$-Value} & \makecell[t]{\textbf{30.0\%}\\\textbf{Threshold}} & \makecell[t]{\textbf{Not Core}\\\textbf{Met}\\\textbf{Threshold}} & \makecell[t]{\textbf{Core Met}\\\textbf{Threshold}} \\ 
\midrule
1 & Intercept & 16074 & 1.22 & 1.90 & 0.0097 & 0.0833 & 0.0000 & - & - & - \\ 
1 & Packages Count & 16074 & 1.22 & 1.90 & 1.1212 & 0.0221 & 0.0000 & \textbf{33.1038} & 3 (0.02\%) & 0 (0.00\%) \\ 
2 & Intercept & 3054 & 2.19 & 2.55 & 0.0325 & 0.1207 & 0.0000 & - & - & - \\ 
2 & Packages Count & 3054 & 2.19 & 2.55 & 1.0393 & 0.0272 & 0.1564 & \textbf{66.8354} & 0 (0.00\%) & 0 (0.00\%) \\ 
5 & Intercept & 1342 & 2.38 & 2.55 & 0.0780 & 0.1338 & 0.0000 & - & - & - \\ 
5 & Packages Count & 1342 & 2.38 & 2.55 & 1.0217 & 0.0338 & 0.5257 & \textbf{79.4868} & 0 (0.00\%) & 0 (0.00\%) \\ 
10 & Intercept & 842 & 2.51 & 2.52 & 0.1329 & 0.1427 & 0.0000 & - & - & - \\ 
10 & Packages Count & 842 & 2.51 & 2.52 & 1.0011 & 0.0376 & 0.9760 & \textbf{1036.2931} & 0 (0.00\%) & 0 (0.00\%) \\ 
20 & Intercept & 230 & 2.29 & 2.64 & 0.2881 & 0.2038 & 0.0000 & - & - & - \\ 
20 & Packages Count & 230 & 2.29 & 2.64 & 1.0462 & 0.0533 & 0.3962 & \textbf{8.7887} & 4 (2.30\%) & 2 (3.57\%) \\ 
\bottomrule
\end{tabular}%
\label{tab:codebase_graph_latex_repository_graph_component_time_series_strategy_logistic_regressionCore_False_packages} 
\end{table}
\end{landscape}
\newpage 

 \newpage 

\begin{landscape}
\begin{table}[H]
\centering
\setlength{\tabcolsep}{2pt}
\caption[Strategy Logistic Regression Core for classes for 1 to 5 and 10]{Strategy Logistic Regression for classes against 1 to 5 and 10 commits to predict Core developers. 
    To identify early indicators of long-term developer engagement, a Logistic Regression model was employed. The model treats the developer category (Sustained vs. Non-Sustained) as a binary dependent variable. }
\begin{tabular}{r l r r r r r r r r r}
\toprule
\makecell[t]{\textbf{First $n$}\\\textbf{Commits}} & \makecell[t]{\textbf{Model}\\\textbf{Parameter}} & \makecell[t]{\textbf{Developers}\\\textbf{\#}} & \makecell[t]{\textbf{Not}\\\textbf{Sustained}\\\textbf{$\mu$}} & \makecell[t]{\textbf{Sustained}\\\textbf{$\mu$}} & \textbf{Odds Ratio} & \textbf{Std. Err} & \textbf{$p$-Value} & \makecell[t]{\textbf{30\%}\\\textbf{Threshold}} & \makecell[t]{\textbf{Not}\\\textbf{Sustained}\\\textbf{Met}\\\textbf{Threshold}} & \makecell[t]{\textbf{Sustained}\\\textbf{Met}\\\textbf{Threshold}} \\ 
\midrule
1 & Intercept & 16074 & 1.26 & 3.92 & 0.0111 & 0.0761 & 0.0000 & - & - & - \\ 
1 & Classes Count & 16074 & 1.26 & 3.92 & 1.0178 & 0.0076 & 0.0209 & \textbf{207.4511} & 1 (0.01\%) & 0 (0.00\%) \\ 
2 & Intercept & 3054 & 2.43 & 6.04 & 0.0344 & 0.1012 & 0.0000 & - & - & - \\ 
2 & Classes Count & 3054 & 2.43 & 6.04 & 1.0102 & 0.0046 & 0.0280 & \textbf{248.8141} & 1 (0.03\%) & 0 (0.00\%) \\ 
5 & Intercept & 1342 & 2.91 & 6.14 & 0.0799 & 0.1050 & 0.0000 & - & - & - \\ 
5 & Classes Count & 1342 & 2.91 & 6.14 & 1.0074 & 0.0044 & 0.0945 & \textbf{227.5778} & 1 (0.08\%) & 0 (0.00\%) \\ 
10 & Intercept & 842 & 3.49 & 6.02 & 0.1304 & 0.1091 & 0.0000 & - & - & - \\ 
10 & Classes Count & 842 & 3.49 & 6.02 & 1.0050 & 0.0042 & 0.2375 & \textbf{239.3208} & 1 (0.13\%) & 0 (0.00\%) \\ 
20 & Intercept & 230 & 4.49 & 7.32 & 0.3172 & 0.1560 & 0.0000 & - & - & - \\ 
20 & Classes Count & 230 & 4.49 & 7.32 & 1.0026 & 0.0044 & 0.5607 & \textbf{117.7201} & 1 (0.57\%) & 1 (1.79\%) \\ 
\bottomrule
\end{tabular}%
\label{tab:codebase_graph_latex_repository_graph_component_time_series_strategy_logistic_regressionCore_False_classes_20} 
\end{table}
\end{landscape}
\newpage 

 \newpage 

\begin{landscape}
\begin{table}[H]
\centering
\setlength{\tabcolsep}{2pt}
\caption[Strategy Logistic Regression Core]{Strategy Logistic Regression for methods against number of commits to predict Core developers}
\begin{tabular}{r l r r r r r r r r r}
\toprule
\makecell[t]{\textbf{Number of}\\\textbf{First}\\\textbf{Commits}} & \makecell[t]{\textbf{Model}\\\textbf{Parameter}} & \makecell[t]{\textbf{Number of}\\\textbf{Developers}} & \makecell[t]{\textbf{Not Core}\\\textbf{$\mu$}} & \textbf{Core $\mu$} & \textbf{Odds Ratio} & \textbf{Std. Err} & \textbf{$p$-Value} & \makecell[t]{\textbf{30.0\%}\\\textbf{Threshold}} & \makecell[t]{\textbf{Not Core}\\\textbf{Met}\\\textbf{Threshold}} & \makecell[t]{\textbf{Core Met}\\\textbf{Threshold}} \\ 
\midrule
1 & Intercept & 16074 & 1.98 & 10.32 & 0.0113 & 0.0749 & 0.0000 & - & - & - \\ 
1 & Methods Count & 16074 & 1.98 & 10.32 & 1.0013 & 0.0006 & 0.0322 & \textbf{2807.8037} & 1 (0.01\%) & 0 (0.00\%) \\ 
1 & Intercept & 16074 & 1.15 & 3.61 & 0.0110 & 0.0759 & 0.0000 & - & - & - \\ 
1 & Methods : packages Count & 16074 & 1.15 & 3.61 & 1.0222 & 0.0064 & 0.0005 & \textbf{166.6467} & 1 (0.01\%) & 1 (0.55\%) \\ 
1 & Intercept & 16074 & 1.94 & 1.86 & 0.0302 & 0.1885 & 0.0000 & - & - & - \\ 
1 & Breadth First & 16074 & 1.94 & 1.86 & 1.0000 & nan & nan & - & - & - \\ 
1 & Depth First & 16074 & 1.94 & 1.86 & 5.2323 & 0.6492 & 0.0108 & - & - & - \\ 
2 & Intercept & 3054 & 6.31 & 17.07 & 0.0353 & 0.0997 & 0.0000 & - & - & - \\ 
2 & Methods Count & 3054 & 6.31 & 17.07 & 1.0008 & 0.0006 & 0.1940 & \textbf{3070.3714} & 1 (0.03\%) & 0 (0.00\%) \\ 
2 & Intercept & 3054 & 1.78 & 4.24 & 0.0342 & 0.1015 & 0.0000 & - & - & - \\ 
2 & Methods : packages Count & 3054 & 1.78 & 4.24 & 1.0154 & 0.0059 & 0.0100 & \textbf{165.0072} & 1 (0.03\%) & 1 (0.95\%) \\ 
2 & Intercept & 3054 & 1.77 & 1.66 & 0.0549 & 0.1645 & 0.0000 & - & - & - \\ 
2 & Breadth First & 3054 & 1.77 & 1.66 & 1.0000 & nan & nan & - & - & - \\ 
2 & Depth First & 3054 & 1.77 & 1.66 & 2.8745 & 0.6427 & 0.1004 & - & - & - \\ 
5 & Intercept & 1342 & 8.73 & 17.05 & 0.0818 & 0.1035 & 0.0000 & - & - & - \\ 
5 & Methods Count & 1342 & 8.73 & 17.05 & 1.0005 & 0.0006 & 0.4395 & \textbf{3301.3251} & 0 (0.00\%) & 0 (0.00\%) \\ 
5 & Intercept & 1342 & 1.32 & 3.35 & 0.0787 & 0.1057 & 0.0000 & - & - & - \\ 
5 & Methods : packages Count & 1342 & 1.32 & 3.35 & 1.0225 & 0.0096 & 0.0198 & \textbf{76.0953} & 1 (0.08\%) & 1 (0.98\%) \\ 
5 & Intercept & 1342 & 1.51 & 1.49 & 0.0906 & 0.1408 & 0.0000 & - & - & - \\ 
5 & Breadth First & 1342 & 1.51 & 1.49 & 1.0000 & nan & nan & - & - & - \\ 
5 & Depth First & 1342 & 1.51 & 1.49 & 6.6218 & 0.7437 & 0.0110 & - & - & - \\ 
10 & Intercept & 842 & 11.75 & 16.96 & 0.1327 & 0.1075 & 0.0000 & - & - & - \\ 
10 & Methods Count & 842 & 11.75 & 16.96 & 1.0003 & 0.0007 & 0.6912 & \textbf{4306.9858} & 0 (0.00\%) & 0 (0.00\%) \\ 
10 & Intercept & 842 & 1.12 & 2.19 & 0.1289 & 0.1098 & 0.0000 & - & - & - \\ 
10 & Methods : packages Count & 842 & 1.12 & 2.19 & 1.0219 & 0.0136 & 0.1099 & \textbf{55.4311} & 1 (0.13\%) & 1 (1.01\%) \\ 
10 & Intercept & 842 & 1.53 & 1.59 & 0.1193 & 0.1632 & 0.0000 & - & - & - \\ 
10 & Breadth First & 842 & 1.53 & 1.59 & 1.0000 & nan & nan & - & - & - \\ 
10 & Depth First & 842 & 1.53 & 1.59 & 4.1905 & 1.2356 & 0.2462 & - & - & - \\ 
20 & Intercept & 230 & 23.05 & 18.77 & 0.3225 & 0.1544 & 0.0000 & - & - & - \\ 
20 & Methods Count & 230 & 23.05 & 18.77 & 0.9999 & 0.0008 & 0.8998 & \textbf{-2978.2890} & 174 (100.00\%) & 56 (100.00\%) \\ 
20 & Intercept & 230 & 1.08 & 2.09 & 0.3160 & 0.1560 & 0.0000 & - & - & - \\ 
20 & Methods : packages Count & 230 & 1.08 & 2.09 & 1.0123 & 0.0162 & 0.4520 & \textbf{25.0150} & 1 (0.57\%) & 2 (3.57\%) \\ 
20 & Intercept & 230 & 1.70 & 1.57 & 0.4717 & 0.2426 & 0.0020 & - & - & - \\ 
20 & Breadth First & 230 & 1.70 & 1.57 & 1.0000 & nan & nan & - & - & - \\ 
20 & Depth First & 230 & 1.70 & 1.57 & 2124278995.8842 & 31654.6738 & 0.9995 & - & - & - \\ 
\bottomrule
\end{tabular}%
\label{tab:codebase_graph_latex_repository_graph_component_time_series_strategy_logistic_regressionCore_False_methods} 
\end{table}
\end{landscape}
\newpage 

 \newpage 


\subsubsection*{Predictive Identification Model}
To determine if early-stage developer behavior can identify future sustained contributors, we employ a Logistic Regression model. The probability $P$ of a developer becoming a \textbf{Sustained Contributor} is modeled using the log-odds (logit) function:

\begin{equation}
    \ln\left(\frac{P}{1-P}\right) = \beta_0 + \beta_1 X_{1} + \beta_2 X_{2} + \dots + \beta_n X_{n}
\end{equation}

Where:
\begin{itemize}
    \item $P$ is the probability that the developer's tenure is ``Sustained''.
    \item $\beta_0$ is the intercept.
    \item $X_n$ represents early-stage metrics (e.g., methods or classes touched in the first 10, 20, or 50 commits).
    \item $\beta_n$ are the coefficients representing the weight of each predictor.
\end{itemize}

The impact of each behavior is quantified using the \textbf{Odds Ratio (OR)}, calculated as:
\begin{equation}
    OR = e^{\beta_n}
\end{equation}
An $OR > 1$ indicates that an increase in the metric (e.g., touching more classes) increases the likelihood of long-term retention. 

\textbf{Model Evaluation}
The predictive power is evaluated using the Area Under the Receiver Operating Characteristic Curve (AUC-ROC). An AUC of 0.5 represents random chance, while an AUC closer to 1.0 indicates high predictive accuracy[cite: 1639].

 \newpage 
